\documentclass[11pt]{article}
\usepackage[utf8]{inputenc}
\usepackage{amsmath,amssymb,amsthm}
\usepackage[margin=1in]{geometry}
\usepackage{hyperref}
\usepackage{xcolor}
\usepackage{enumitem}
\usepackage{newunicodechar}
\newunicodechar{ℝ}{\ensuremath{\mathbb{R}}}
\newunicodechar{⟨}{\ensuremath{\langle}}
\newunicodechar{⟩}{\ensuremath{\rangle}}
\newunicodechar{λ}{\ensuremath{\lambda}}
\newunicodechar{α}{\ensuremath{\alpha}}
\newunicodechar{ν}{\ensuremath{\nu}}
\newunicodechar{ρ}{\ensuremath{\rho}}
\newunicodechar{χ}{\ensuremath{\chi}}
\newunicodechar{²}{\ensuremath{^2}}
\newunicodechar{₀}{\ensuremath{_0}}
\newunicodechar{₁}{\ensuremath{_1}}
\newunicodechar{₂}{\ensuremath{_2}}
\newunicodechar{⁻}{\ensuremath{^{-}}}
\newunicodechar{→}{\ensuremath{\to}}
\newunicodechar{↦}{\ensuremath{\mapsto}}

\newcommand{\userbox}[1]{\par\medskip\begin{quote}\small\textbf{\textcolor{blue!60}{DM:}} #1\end{quote}\medskip}
\newcommand{\claudebox}[1]{\par\medskip\begin{quote}\small\textbf{\textcolor{orange!60!black}{Claude:}} #1\end{quote}\medskip}
\newcommand{\gptbox}[1]{\par\medskip\begin{quote}\small\textbf{\textcolor{green!50!black}{GPT-5.6 Sol:}} #1\end{quote}\medskip}

\newcommand{\tideref}[2][]{\href{https://therisensea.org/tides/#2/}{\ifx&#1&\texttt{#2}\else#1\fi}}
\newcommand{\experimentref}[2][]{\href{https://therisensea.org/experiments/#2/}{\ifx&#1&\texttt{#2}\else#1\fi}}
\newcommand{\reviewref}[2][]{\href{https://therisensea.org/reviews/#2/}{\ifx&#1&\texttt{#2}\else#1\fi}}

\newcommand{\Cov}{\operatorname{Cov}}
\newcommand{\Var}{\operatorname{Var}}
\emergencystretch=3em

\title{Review retrospective:\\\emph{the germbij forward completion and the class-(c) arc}}
\author{Claude (Anthropic), with GPT-5.6 Sol consults}
\date{2026-08-10}

\begin{document}
\maketitle

\begin{abstract}
The second review of the germbij formalisation
arc\footnote{\reviewref[The first]{2026-08-10-01-29-review-germbij-new-content}
covered PRs \#93--\#103.} covers the seven files that landed after it:
the forward-programme completion (PRs \#104--\#110 --- the numerator
integration, the outer tails, the generic asymptotic division, and the
public \texttt{rescaledMoment\_hasExpansion}) and the class-(c)
semi-quasi-homogeneous arc (PRs \#113--\#116 --- the exact
quasi-homogeneous moment laws, the one-grade covariance engine, and
grade recovery). The independent reviewer found \emph{no mathematical
fidelity defects in either arc}: the division recursion, the tail
estimates, the exponent conventions ($-\langle q, \alpha +
\mathbf{1}\rangle$ against the note's $\ell(\alpha)$), the sign of the
covariance limit, and the constant-killing equivalence all verified,
independently agreeing with the session's own hand checks. What it
found instead were three boundary-honesty findings --- places where a
docstring or a natural reading claims slightly more than the theorems
deliver --- each with a documentation-level fix that the review
implemented under the statement-invariance gates (fingerprint
byte-identical, build green). Twenty-one proof-refactor suggestions are
recorded for a dedicated simplification pass.
\end{abstract}

\section{Setting}

The germbij fidelity-review arc resumed at the user's direction after
the perimeter tides closed the first review's findings. Since that
review's baseline, two programmes had landed on \texttt{main}: the
forward-expansion programme's completion --- stages 5c-ii through 7,
proving that rescaled posterior moments of continuous polynomial-growth
observables are order-$N$ asymptotic polynomials at $0^+$ --- and the
scoping-consulted fragment of the note's \S7.4 item (c). Neither had
been independently reviewed. The review ran as one pass with the
consult split per arc, against the note's \S3 (forward direction),
\S7.4(c) (quasi-homogeneous recovery), the two archived architecture
consults, and the upstream definition corpus. The note's two
\verb|\leanref| anchors into \texttt{QhMomentRecovery} were verified to
point at the correct lines before the consult.

\section{What was reviewed}

\textbf{Forward completion.} \texttt{Multi/NumeratorExpansion.lean}
(the DCT window piece with coefficients $a_j = \int P\, e^{-T_2}\,
P_j$), \texttt{Multi/NumeratorTails.lean} (both outer tails beyond all
orders; the packaging \texttt{numerator\_hasExpansion}),
\texttt{Multi/AsymptoticDivision.lean} (the generic division lemma with
the recursion $c_j = (a_j - \sum_{i=1}^{j} b_i c_{j-i})/b_0$), and
\texttt{Multi/ForwardTheorems.lean} (the public face:
\texttt{rescaledMoment\_hasExpansion}, its existence form, the monomial
instances, and the comparison face
\texttt{momentCoeff\_eq\_of\_isLittleO}).

\textbf{Class (c).} \texttt{Multi/QhMomentRecovery.lean} (exact
unnormalized and normalized monomial moment laws for a
quasi-homogeneous potential; the moment-ratio recovery
$t^{\langle q,\alpha\rangle}\langle x^\alpha\rangle_t = M_\alpha/M_0$,
all multi-indices including odd; and the single-temperature
comparison form), the one-grade file
\texttt{Multi/GradeComparison.lean} (the pointwise engine
$|e^{-W}-1+W| \leq W^2$ and the one-grade difference limit
$(M_2(h)-M_1(h))/h^\rho \to -\Cov_P(A,Q)$), and
\texttt{Multi/GradeRecovery.lean} (covariance injectivity and
\texttt{grade\_eq\_of\_normalized\_rates}).

\section{Fidelity findings}

\textbf{No mathematical defects in either arc.} The reviewer
re-derived the historically error-prone pieces: the division
recursion's indexing and its handling of negative $b_0$ (through
$|g| \geq |b_0|/2$ eventually), the eventual decomposition of the
numerator into window plus tails with the domain indicator collapsing
on the window, the integrability supplied to the positivity route for
$a_0(1) = \int e^{-T_2} > 0$, the exponent bookkeeping
$-(\sum_i q_i + \sum_i q_i\alpha_i) = -\ell(\alpha)$ with the volume
factor cancelling in the normalized law, the sign of the covariance
limit under $V_2 - V_1 \to Q$, and the equivalence of the Lean's
$Q(0) = 0$ constant-killing with the note's positive-weighted-degree
argument. All checks converged with the session's independent
pre-triage, recorded in the review log before the responses arrived.

Three boundary-honesty findings, all of the ``narrower than the
prose'' genre:

\begin{enumerate}[leftmargin=2em]
\item \textbf{Fixed rescaled observables (forward).} The expansion
theorems hold for a fixed observable of the \emph{rescaled} variable.
The note's fixed original-scale $\varphi(w)$ becomes the $q$-dependent
$z \mapsto \varphi(qz)$ after rescaling, and the final
observable-Taylor packaging is not in these files; the monomial
instances (where $\varphi_m(qz) = q^k\cdot\texttt{monomialTest}$)
carry the content the recovery theorems consume. The top docstring's
``the germbij note's forward direction'' read broader than that.
\item \textbf{Global exact laws vs the localized cutoff law (SQH).}
The note states $\int x^\alpha e^{-tP}\chi = t^{-\ell(\alpha)}M_\alpha
+ O(e^{-ct})$; the Lean proves the \emph{global} model exactly, with
no cutoff and no error term --- a faithful strengthening for that
model, but the localized bridge is unformalised. Moreover the exact
identities are junk-value-consistent, taking no integrability or
positivity hypotheses, so ``recovery'' is substantive only when the
reference moments are finite with $M_0 > 0$ (the note's standing
assumptions, not taken in the file).
\item \textbf{The generic engine vs the SQH specialization.}
\texttt{GradeComparison} is the correct one-grade engine, but nothing
yet constructs the correction families from weighted-homogeneous
pieces, identifies the rate exponent $\rho \in \mathbb{N}$ with the
cleared grade difference, or runs the finite grade induction ---
so \texttt{GradeRecovery}'s ``reaching the designated endpoint'' was
optimistic; ``the conditional one-grade endpoint'' is accurate.
\end{enumerate}

\section{Simplifications applied}

The three documentation-level fixes the reviewer itself recommended
for the findings above --- and nothing else. The
\texttt{ForwardTheorems} module docstring now says ``the
rescaled-moment core of the germbij forward direction'' and names the
two non-goals (observable-Taylor packaging, Wick evaluation of the
coefficient integrals); the monomial-instance docstring states the
original-scale bridge $\varphi_m(qz) = q^k\cdot\texttt{monomialTest}$
explicitly; \texttt{QhMomentRecovery} carries the two boundary notes
(junk-value consistency and the unformalised localized law); and
\texttt{GradeRecovery}'s endpoint claim is qualified. Statement
invariance was checked mechanically (the public-signature fingerprint
of all seven files is byte-identical before and after) and the build
is green. The affirmation pass endorsed all three changes. The 21
proof-refactor suggestions (growth-normalization and Gaussian
domination helpers, integrand abbreviations, the reusable
eventual-denominator lemma, finite-sum linearity helpers, and the
rest) are recorded verbatim in the review log as a batch for a
dedicated simplification pass.

\section{Disagreements and open questions}

No Claude/GPT disagreements; the consults' findings coincide with the
pre-triage. Escalated to the human, as candidate future work rather
than defects:
\begin{itemize}[nosep]
\item the localized cutoff law (the $O(e^{-ct})$ bridge from the
global model to the note's $\chi$-cut statement) --- one tide, the
same loss-gap mechanism as the inverse-side cutoff removal;
\item the SQH instantiation of the one-grade engine
(weighted-homogeneous corrections, $\rho$ = cleared grade difference,
finite induction) --- the consult-designated continuation if class (c)
is ever re-prioritised;
\item the observable-Taylor packaging for fixed original-scale
observables, if the note's forward statement is wanted verbatim.
\end{itemize}

\section{Lessons}

The pattern from the first review repeats cleanly: on this seabed the
failure mode is never a wrong statement --- it is prose (docstrings,
completion banners) claiming the neighbourhood of what was proved
rather than exactly what was proved. Fidelity reviews should diff the
hypothesis list and the quantifier structure against the claim
sentence by default; both consults were most valuable exactly there.
Also worth keeping: verifying the note's \verb|\leanref| line numbers
against the files takes seconds and closes a whole class of
anchor-drift findings before the consult runs.

\section{Follow-ups}

\begin{itemize}[nosep]
\item The recorded 21-item simplification batch (review log,
``Triaged findings'').
\item The three candidate tides above, at the user's discretion.
\item Review the perimeter files (HessianBridge, ShiftNormalization,
CutoffRemoval) once PRs \#109/\#112/\#115 merge — they are unmerged
and were deliberately out of scope here.
\end{itemize}

\end{document}
